% Figure: external-compression supersonic inlet shock system. Two
% compression ramps throw two oblique shocks that meet at the cowl lip;
% a terminal normal shock just inside the lip drops the flow to
% subsonic. Schematic, not to scale; angles chosen for legibility.
\begin{figure}[htbp]
\centering
\begin{tikzpicture}[scale=1.0]
% ramps (the compression surface)
% first ramp from (0,0) to (3,0.9), second ramp to (5,2.1)
\draw[engblue, very thick] (0,0) -- (3,0.9) -- (5,2.1);
\node[engblue, font=\scriptsize, anchor=north] at (1.5,0.25) {ramp 1};
\node[engblue, font=\scriptsize, anchor=north] at (4.1,1.35) {ramp 2};
% cowl lip (upper body)
\draw[engblue, very thick] (5.2,3.0) -- (6.6,3.0);
\node[engblue, font=\scriptsize, anchor=south] at (5.9,3.05) {cowl};
% engine-face duct (subsonic)
\draw[enggray, thick] (5.0,2.1) -- (6.6,2.4);
\draw[enggray, thick] (6.6,3.0) -- (6.6,2.4);
\node[enggray, font=\scriptsize] at (6.95,2.7) {engine};
\node[enggray, font=\scriptsize] at (6.95,2.45) {face};
% freestream arrow
\draw[engblack, -{Latex[length=2mm]}, thick] (-1.4,1.6) -- (-0.3,1.6);
\node[engblack, font=\scriptsize, anchor=south] at (-0.85,1.62) {$M_\infty=3$};
% first oblique shock: from ramp-1 corner (0,0) up to lip
\draw[engred, thick] (0,0) -- (5.2,3.0);
\node[engred, font=\scriptsize, anchor=west, rotate=30] at (2.4,1.5) {shock 1};
% second oblique shock: from ramp-2 corner (3,0.9) up to lip
\draw[engamber, thick] (3,0.9) -- (5.2,3.0);
\node[engamber, font=\scriptsize, anchor=north west] at (3.7,1.6) {shock 2};
% terminal normal shock just inside the lip
\draw[englime, very thick] (5.2,2.1) -- (5.2,3.0);
\node[englime, font=\scriptsize, anchor=west] at (5.25,2.45) {normal};
\node[englime, font=\scriptsize, anchor=west] at (5.25,2.2) {shock};
% Mach-number labels in each zone
\node[engblack, font=\scriptsize] at (1.3,1.05) {$M_1=3$};
\node[engblack, font=\scriptsize] at (3.7,1.05) {$M_2$};
\node[engblack, font=\scriptsize] at (4.7,1.45) {$M_3$};
\node[engblack, font=\scriptsize] at (5.9,2.6) {$M_4<1$};
\end{tikzpicture}
\caption[External-compression supersonic inlet]{The shock system of a
two-ramp external-compression inlet at $M_\infty=3$, drawn
schematically. The freestream meets the first ramp and turns through an
oblique shock (shock~1), then the second ramp and a second oblique
shock (shock~2); the two ramps are angled so both shocks focus on the
cowl lip, the shock-on-lip condition that captures the full stream tube
without spillage. Each oblique shock is weak, so each loses little
stagnation pressure. A terminal normal shock just inside the lip drops
the now-decelerated supersonic flow to the subsonic state the engine
face requires. Stacking two weak oblique shocks ahead of a weak
terminal normal shock recovers far more stagnation pressure than a
single strong normal shock would at $M_\infty=3$, the quantitative
content of the inlet case study.}
\label{fig:vol03ch13:inlet-shocks}
\end{figure}
